\begin{examplebox}
\textbf{例 1（由存在的角度反求离心率）}\quad
椭圆上存在点 $P$ 满足 $\angle F_1PF_2=120^\circ$，求 $e$ 的范围。

\textbf{解：}存在这样的点，当且仅当最大张角不少于 $120^\circ$：
\[
e\ge\sin\frac{120^\circ}{2}
=\sin60^\circ=\frac{\sqrt3}{2}.
\]
结合椭圆 $0<e<1$，得
\[
\boxed{e\in\left[\frac{\sqrt3}{2},1\right).}
\]
\textbf{点拨：}“存在”允许在短轴端点取到等号，所以下界要包含。

\medskip
\textbf{例 2（由离心率求最大张角）}\quad
已知椭圆离心率 $e=\frac12$，求 $\theta$ 的最大值。

\textbf{解：}
\[
\sin\frac{\theta_{\max}}2=\frac12
\quad\Longrightarrow\quad
\frac{\theta_{\max}}2=30^\circ
\quad\Longrightarrow\quad
\boxed{\theta_{\max}=60^\circ.}
\]
\textbf{点拨：}在短轴端点才能取到最大值。

\medskip
\textbf{例 3（结合具体椭圆验证“短轴最大”）}\quad
设 $\frac{x^2}{25}+\frac{y^2}{16}=1$，比较短轴端点与点
$P=(3,\frac{16}{5})$ 的焦点张角。

\textbf{解：}这里 $a=5,b=4,c=3,e=\frac35$。
短轴端点的张角满足
\[
\sin\frac{\theta_{\max}}2=\frac35.
\]
对于 $P=(3,\frac{16}{5})$，它确实在椭圆上，且
\[
r_1=a+ex=5+\frac35\cdot3=\frac{34}{5},
\qquad
r_2=a-ex=\frac{16}{5}.
\]
所以
\[
\sin^2\frac\theta2
=1-\frac{16}{(34/5)(16/5)}
=\frac9{34},
\qquad
\sin\frac\theta2=\frac3{\sqrt{34}}<\frac35.
\]
故该点的张角严格小于短轴端点的最大张角。
\textbf{点拨：}离开 $x=0$ 后，两焦半径一长一短，张角会变小。

\medskip
\textbf{例 4（直角、钝角的存在性，一题两问）}\quad
分别求椭圆上存在焦点张角为直角、存在焦点张角为钝角时 $e$ 的条件。

\textbf{解：}
\begin{enumerate}
    \item \textbf{存在直角：}
    $\theta_{\max}\ge90^\circ$，因此
    \[
    \boxed{e\ge\sin45^\circ=\frac{\sqrt2}{2}.}
    \]
    \item \textbf{存在钝角：}
    $\theta_{\max}>90^\circ$，因此
    \[
    \boxed{e>\frac{\sqrt2}{2}.}
    \]
\end{enumerate}
两题均还要满足 $e<1$。
\textbf{点拨：}“直角”可取边界，“钝角”必须严格超过边界。

\medskip
\textbf{例 5（变式：固定角度反求动点坐标）}\quad
在椭圆 $\frac{x^2}{25}+\frac{y^2}{9}=1$ 上，求满足
$\angle F_1PF_2=60^\circ$ 的所有点 $P$。

\textbf{解：}由 $a=5,b=3,c=4,e=\frac45$ 及焦半径形式，
\[
\sin^2\frac\theta2
=1-\frac{b^2}{a^2-e^2x^2}.
\]
代入 $\sin^230^\circ=\frac14$，有
\[
\frac14=1-\frac{9}{25-\frac{16}{25}x^2}
\quad\Longrightarrow\quad
25-\frac{16}{25}x^2=12
\quad\Longrightarrow\quad
x^2=\frac{325}{16}.
\]
再代入椭圆方程，
\[
y^2=9\left(1-\frac{x^2}{25}\right)=\frac{27}{16}.
\]
所以共有四个点：
\[
\boxed{P=\left(\pm\frac{5\sqrt{13}}4,\ \pm\frac{3\sqrt3}4\right),}
\]
两个正负号可以独立选取。
\textbf{点拨：}本例展示了“不等式判断存在性”之后，还能用焦半径公式精确定位动点。
\end{examplebox}
